\pdfmapfile{+symbols.map}
\pdfmapfile{+cm.map}
\pdfmapfile{+cmextra.map}
\pdfmapfile{+lm.map}
\documentclass[11pt,letterpaper]{article}
\usepackage[margin=1in]{geometry}
\usepackage{amsmath,amssymb,amsthm}
\usepackage[T1]{fontenc}
\usepackage{lmodern,microtype}
\usepackage[colorlinks=true,linkcolor=blue,urlcolor=blue,citecolor=blue]{hyperref}
\newtheorem{theorem}{Theorem}
\newcommand{\Q}{\mathbb Q}
\newcommand{\R}{\mathbb R}
\newcommand{\Syl}{\operatorname{Syl}}
\newcommand{\Sres}{\operatorname{Sres}}
\title{Disproof of Conjecture 00000005613}
\author{}
\date{20 October 2026}
\begin{document}
\maketitle
\vspace{-1.5em}
\begin{abstract}
A third-order subresultant coefficient minor can vanish while its
corresponding principal minor is nonzero. For the monic, squarefree,
coprime, strictly interlacing polynomials $f(x)=x^3-4x$ and
$g(x)=x^2-1$, the first subresultant is $-3x$: its constant-coefficient
minor is zero, whereas its principal coefficient is $-3$. The example
uses the actual Sylvester submatrix construction and is verified in
Lean 4 using only the standard library.
\end{abstract}

\section{The precise clause and standard convention}
Conjecture 00000005613\footnote{\href{https://github.com/The-Last-Math-Competition/The-Last-Math-Competition/blob/main/conjectures/00000005613.md}{Official conjecture 00000005613 on GitHub}, checked 20 October 2026.}
asserts that ``a minor is zero if and only if the corresponding
principal minor is zero'' in a discussion of third-order subresultant
minors. We disprove this clause using standard subresultant coefficient
minors and their canonical leading principal minor.

For degrees $p,q$ and $i<\min(p,q)$, the $i$-th Sylvester submatrix
has rows given by
\[
 x^{q-i-1}f,\ldots,f,\quad x^{p-i-1}g,\ldots,g,
\]
written as coefficient vectors in descending degree. Its dimensions
are $(p+q-2i)\times(p+q-i)$. For a matrix $A$, let $M_{k,t}(A)$ be
the determinant using its first $k$ rows, first $k-1$ columns, and
column $k+t$ (one-based indexing). Thus $M_{k,0}$ is the leading
principal minor. The standard polynomial subresultant is
\[
 \Sres_i(f,g)(x)=\sum_{j=0}^{i}
 M_{p+q-2i,i-j}\bigl(\Syl_i(f,g)\bigr)x^j;
\]
its principal coefficient is $M_{p+q-2i,0}$.
These conventions are stated explicitly in Kerber,
Definitions 2.3.10--2.3.12\cite{Kerber}.

The printed conjecture does not separately define ``spectrum,''
``conversion,'' or ``corresponding.'' This disproof uses the standard
coefficient/principal correspondence above. It makes no claim about an
unspecified alternative correspondence or about the other undefined
spectral assertions. In particular, the example is not merely an
arbitrary matrix with unrelated minors. The target is an \emph{individual}
coefficient minor, not the vanishing of the whole subresultant polynomial.
In this example the whole first subresultant and its principal coefficient
are both nonzero; that different equivalence is not being contradicted.

\newpage
\section{The counterexample}
Take
\[
 f(x)=x^3-4x=(x+2)x(x-2),\qquad
 g(x)=x^2-1=(x+1)(x-1).
\]
Their roots strictly alternate:
\[
 -2<-1<0<1<2.
\]
Both polynomials are monic and have distinct real roots. The
coefficient-level identity
\begin{equation}\label{eq:bezout}
 x f(x)+(3-x^2)g(x)=-3
\end{equation}
proves they are coprime over $\Q$, $\R$, or $\mathbb C$: a common
polynomial divisor would divide the nonzero constant $-3$.

Here $p=3$, $q=2$, $i=1$, so the relevant minors have order
$p+q-2i=3$. The rows of the first Sylvester submatrix are
$f,xg,g$, in the basis $x^3,x^2,x,1$:
\[
 S=\Syl_1(f,g)=
 \begin{pmatrix}
 1&0&-4&0\\
 1&0&-1&0\\
 0&1&0&-1
 \end{pmatrix}.
\]
This matrix is generated from the polynomial coefficients and the
prescribed shifts. Its principal coefficient is
\begin{equation}\label{eq:principal}
 M_{3,0}(S)=
 \det\begin{pmatrix}
 1&0&-4\\1&0&-1\\0&1&0
 \end{pmatrix}=-3\ne0.
\end{equation}
The corresponding constant-coefficient minor is
\begin{equation}\label{eq:constant}
 M_{3,1}(S)=
 \det\begin{pmatrix}
 1&0&0\\1&0&0\\0&1&-1
 \end{pmatrix}=0.
\end{equation}
The first determinant equals $-3$ by expansion along its second
column. The second is zero because its first two rows agree.
Consequently
\[
 \Sres_1(f,g)(x)=M_{3,0}(S)x+M_{3,1}(S)=-3x.
\]

\begin{theorem}
The vanishing-equivalence clause of Conjecture 00000005613 is false
under the standard subresultant coefficient/principal convention,
even for monic squarefree coprime polynomials with strictly interlacing
real roots.
\end{theorem}
\begin{proof}
Equations \eqref{eq:constant} and \eqref{eq:principal} give a zero
third-order coefficient minor whose corresponding principal minor is
nonzero. This directly contradicts the asserted equivalence.
\end{proof}

\newpage
\section{Formal verification and its boundary}
\texttt{Subresultant.lean} defines polynomials by ascending rational
coefficients, and generates the Sylvester submatrix for arbitrary
$p,q,i$ by shifted coefficients. Its determinant is ordinary recursive
Laplace expansion. The coefficient-minor and subresultant definitions
implement the convention in Section 1 directly.

The file proves exact degrees and monicity; the two determinant values;
the first subresultant formula; the displayed matrix from its generated
entries; rational factorizations and exact rational root sets; strict
root ordering; and the B\'ezout identity. It also verifies the two
factorizations and \eqref{eq:bezout} by finite coefficient-list
multiplication, independently of evaluating $x$.

The final theorem is \texttt{conjecture\_00000005613\_false}.
It negates \texttt{VanishingEquivalence}, the claimed equivalence restricted to
monic degree-three/degree-two inputs satisfying the explicit
factorization-based strict-interlacing condition. This is a restricted
instance of the stated universal clause, not an assertion that every
remaining undefined term has been formalized.

\textbf{Scalar-domain boundary.} Lean's scalar type here is
\texttt{Rat}, not a formal type of real or complex numbers. The
coefficient identities are checked exactly over $\Q$; their transfer
along the usual embeddings into $\R$ and $\mathbb C$, and the resulting
real-root and polynomial-coprimality interpretations, are the elementary
mathematical steps explained above. The theorem
\texttt{no\_common\_root} itself quantifies over rational inputs; it is
not advertised as a formal general polynomial gcd theorem.

Strict compilation under Lean 4.31.0 succeeds with only \texttt{Std}.
All theorem axiom closures use at most \texttt{propext},
\texttt{Classical.choice}, and \texttt{Quot.sound}; there are no custom
axioms, admitted proofs, or native-decision shortcuts.
The supplied README gives the strict build and kernel-replay commands.
The separate Python script recomputes the generic matrix, determinants,
coefficient products, and B\'ezout identity using exact rational
arithmetic; it is supplementary, not a substitute for the Lean proof.

\begin{thebibliography}{1}
\bibitem{Kerber}
Michael Kerber.
\emph{Geometric Algorithms for Algebraic Curves and Surfaces}.
Dissertation, Universit\"at des Saarlandes, 2009,
Section 2.3, Definitions 2.3.10--2.3.12, p. 30.
\href{https://www.geometrie.tugraz.at/kerber/kerber_papers/kerber_diss.pdf}{Author-hosted full text}.
\end{thebibliography}
\end{document}
